% chapters/10/contoh.tex
\section{Dua Contoh Analisis}
\label{sec:contoh-bab-10}

\subsection{Contoh Pertama: Menghitung Biaya, Perkiraan, dan Jumlahnya}

\begin{contohbox}
Sebuah peta bersih berukuran lima kali empat petak, tanpa dinding. Titik
awalnya berada pada petak $(0, 3)$ dan titik tujuannya pada petak $(4, 0)$.
Heuristik yang dipakai adalah jarak Manhattan.

Jarak Manhattan antara titik awal dan titik tujuan adalah

\[
h_{\text{awal}} = |4 - 0| + |0 - 3| = 7 .
\]

Tabel~\ref{tab:hitung-fgh} memuat perhitungan $g$, $h$, dan $f$ untuk beberapa
petak.
\end{contohbox}

\begin{table}[htbp]
\centering
\caption{Perhitungan biaya, perkiraan, dan jumlahnya pada beberapa petak}
\label{tab:hitung-fgh}
{\setlength{\tabcolsep}{4pt}\footnotesize%
\begin{tabularx}{\textwidth}{@{}C{2.2cm} C{2.4cm} C{2.4cm} C{2.0cm} Y@{}}
\toprule
\headrow \textbf{Petak} & \textbf{$g$ (langkah dari awal)} & \textbf{$h$ (perkiraan ke tujuan)} & \textbf{$f = g + h$} & \textbf{Keterangan} \\
\midrule
$(0,3)$ & 0 & 7 & 7 & Titik awal. \\
\addlinespace
$(1,3)$ & 1 & 6 & 7 & Satu langkah ke kanan. \\
\addlinespace
$(2,2)$ & 3 & 4 & 7 & Dua langkah menuju tujuan, satu langkah menjauh. \\
\addlinespace
$(3,1)$ & 5 & 2 & 7 & Lima langkah ditempuh, dua langkah tersisa. \\
\addlinespace
$(4,0)$ & 7 & 0 & 7 & Titik tujuan. \\
\bottomrule
\end{tabularx}}
\tabsource{Nilai $f$ sama pada seluruh petak, yaitu 7. Sifat
itu berlaku pada peta tanpa dinding, sebab setiap langkah mengurangi perkiraan
sisa jarak tepat satu.}
\end{table}

Baris ketiga menarik untuk diperhatikan. Petak $(2,2)$ menuntut tiga langkah,
yaitu satu langkah lebih banyak daripada jarak mendatarnya dari titik awal.
Meskipun demikian, nilai $h$-nya berkurang sebanyak itu pula, sehingga nilai
$f$-nya tetap 7. Itulah sebabnya petak itu tetap dianggap sama menjanjikannya
dengan petak-petak lain pada jalur terpendek.

Sifat itu menjelaskan mengapa A* pada peta tanpa dinding memeriksa begitu
banyak petak: seluruh petak memiliki nilai $f$ yang sama, sehingga heuristiknya
tidak memberi petunjuk arah sama sekali. Manfaat heuristik baru terasa ketika
ada dinding: petak yang terpaksa memutar memperoleh $g$ yang lebih besar tanpa
pengurangan $h$ yang sepadan, sehingga nilai $f$-nya membesar dan A* menunda
pemeriksaannya.

\subsection{Contoh Kedua: Menganalisis Hasil Perbandingan}

\begin{contohbox}
Pada peta bab ini, ketiga heuristik menghasilkan biaya jalur yang sama, yaitu
50, dengan panjang 51 petak. Banyaknya petak yang diperiksa berbeda-beda:
Manhattan 274 petak, Euclidean 353 petak, dan tanpa perkiraan 437 petak.
\end{contohbox}

Perbandingan itu dapat dinyatakan dalam bentuk rasio, sehingga lebih mudah
ditafsirkan.

\begin{itemize}
  \item Euclidean memeriksa $353 \div 274 \approx 1{,}29$ kali lipat petak
        Manhattan, yaitu sekitar 29 persen lebih banyak.
  \item Tanpa perkiraan memeriksa $437 \div 274 \approx 1{,}59$ kali lipat,
        yaitu sekitar 59 persen lebih banyak.
\end{itemize}

Biaya jalurnya tetap sama pada ketiganya. Inilah yang dimaksud
dengan pengaruh heuristik terhadap \textbf{kinerja} pencarian, bukan terhadap
\textbf{mutu} hasilnya: yang berubah hanyalah banyaknya pekerjaan yang perlu
dilakukan. Pada peta yang kecil, pekerjaan tambahan itu tidak terasa. Pada peta
yang luas, ia menentukan apakah permainan masih dapat berjalan pada laju
bingkai yang wajar.

Perlu dicatat pula bahwa rasio itu \textbf{tidak tetap} untuk peta yang
berbeda. Pada peta yang penuh liku, keunggulan Manhattan menjadi jauh lebih
besar, sebab Dijkstra terpaksa memeriksa hampir seluruh petak yang dapat
dilalui. Sebaliknya, pada peta yang jalurnya hampir lurus, ketiganya hampir
sama saja. Karena itu, perbandingan semacam ini sebaiknya dilakukan pada peta
yang benar-benar dipakai, bukan pada peta contoh.

\begin{catatabox}
Satu hal yang tidak terlihat pada perbandingan itu adalah \textbf{mutu jalur}
apabila heuristiknya tidak memenuhi syarat. Selama heuristiknya tidak
melebihkan jarak, seluruh baris pada tabel akan memiliki biaya yang sama. Begitu
heuristiknya melebihkan jarak---misalnya karena dikalikan dua---biaya jalurnya
dapat membesar, dan itulah satu-satunya tanda bahwa ada yang keliru. Karena itu,
pemeriksaan biaya jalur harus selalu dijalankan bersama perbandingan waktu
pencarian.
\end{catatabox}
